\section{Capability Accumulation and Graduation}
\label{sec:graduation}

Activation is not development by itself. A supported relation can operate for years while remaining exactly as dependent on support as it was at entry. The state variable $q_{e,t}$ is introduced to separate operation from durable learning.

During operation, capability follows equation~\eqref{eq:capability-law}. The term $\Lambda_e(\Omega_t)$ should be interpreted broadly as capability conversion, not merely technical learning. It can include learning-by-doing, quality improvement, managerial routines, supplier development, standards compliance, export competence, absorptive capacity, technology adaptation, and other durable changes that improve future unsupported viability. The term $\delta_e$ captures capability erosion through obsolescence, organizational destruction, supplier loss, worker exit, repeated macroeconomic disruption, or other mechanisms that destroy accumulated productive state.

Unsupported viability requires $q_{e,t}\ge q_e^*(\Omega_t)$. The threshold is reduced form: it can summarize the capability required to operate without the extraordinary policy wedge given the current infrastructure, market access, input structure, and other conditions.

\subsection{Baseline constant-parameter case}

Consider one relation operated continuously during a support interval:
\begin{equation}
q_{t+1}=(1-\delta)q_t+\lambda,
\qquad
0<\delta<1,
\quad \lambda>0,
\quad q_0<q^*.
\label{eq:baseline}
\end{equation}
The solution is
\begin{equation}
q_t
=(1-\delta)^tq_0
+\frac{\lambda}{\delta}
\left[1-(1-\delta)^t\right].
\label{eq:closed-form}
\end{equation}
Hence
\begin{equation}
q_\infty=\frac{\lambda}{\delta}.
\end{equation}

\begin{proposition}[Graduation feasibility]
\label{prop:graduation}
Let $0<\delta<1$ and $q_0<q^*$. If $\lambda/\delta\le q^*$, the relation never reaches $q^*$ in finite time. If $\lambda/\delta>q^*$, a finite graduation time exists.
\end{proposition}

\begin{proof}
Equation~\eqref{eq:closed-form} writes $q_t$ as a convex combination of $q_0$ and $q_\infty=\lambda/\delta$. If $q_\infty<q^*$, both endpoints are below $q^*$, so every finite $q_t<q^*$. If $q_\infty=q^*$, then
\[
q_t=q^*-(1-\delta)^t(q^*-q_0)<q^*
\]
for every finite $t$. If $q_\infty>q^*$, the sequence converges monotonically from $q_0$ to a level strictly above the threshold, so a finite integer crossing time exists.
\end{proof}

The knife-edge case $\delta=1$ is separate. Then $q_{t+1}=\lambda$ during operation, so graduation occurs after one period if and only if $\lambda\ge q^*$.

\begin{corollary}[Minimum support duration]
\label{cor:duration}
Suppose $0<\delta<1$ and
\[
q_0<q^*<\frac{\lambda}{\delta}.
\]
The minimum integer number of operating periods required to reach unsupported viability is
\begin{equation}
k^*
=
\left\lceil
\frac{
\ln\left(
\dfrac{\lambda/\delta-q^*}{\lambda/\delta-q_0}
\right)
}{
\ln(1-\delta)
}
\right\rceil.
\label{eq:min-duration}
\end{equation}
\end{corollary}

\begin{proof}
Rewrite equation~\eqref{eq:closed-form} as
\[
q_t
=
\frac{\lambda}{\delta}
-(1-\delta)^t
\left(
\frac{\lambda}{\delta}-q_0
\right).
\]
The condition $q_t\ge q^*$ is equivalent to
\[
(1-\delta)^t
\le
\frac{\lambda/\delta-q^*}{\lambda/\delta-q_0}.
\]
Both sides lie in $(0,1)$. Taking logs and dividing by the negative number $\ln(1-\delta)$ yields the stated lower bound on $t$; taking the ceiling gives the minimum integer duration.
\end{proof}

Define the reduced-form graduation margin
\begin{equation}
\gamma
=
\frac{\lambda}{\delta q^*}.
\label{eq:gamma}
\end{equation}
For $q_0<q^*$ and $0<\delta<1$, $\gamma>1$ is exactly the condition for eventual finite graduation in the constant-parameter case. The margin is useful as a diagnostic, but it is not a welfare measure and it becomes incomplete when the environment changes endogenously.

The substantive interpretation is simple. A policy can keep an activity alive without changing the state that determines unsupported viability. Lengthening support does not solve that problem. The developmental question is whether operation under support changes $\lambda$, $\delta$, $q^*$, or the accumulated state $q$ sufficiently to create a path out of support.
